% s-382-p382-m001; original PDF 382
\PoemStart{s-382-p382-m001}{Rath eines Alten.}
\begin{BookVerse}
\VerseLine{0}{An solche Weiblein nicht dein Wünschen hefte,}
\VerseLine{1}{Die fleischgeword’ner „guter Wille“ sind!}
\VerseLine{1}{So hülflos rührend ist solch gutes Kind,}
\VerseLine{0}{Das an das Grösste wagt die kleinen Kräfte.}
\Stanza
\VerseLine{1}{Sieh nur nicht hin, Mann! stell’ dich lieber blind!}
\VerseLine{0}{Sonst kommt das Mitleid, das die Stärksten äffte:}
\VerseLine{1}{Man springt hinzu, man rettet, hilft geschwind}
\VerseLine{0}{Und mengt sich in verfahrene Geschäfte.}
\Stanza
\VerseLine{1}{Das zarte Füsschen, seidenweich beschuht,}
\VerseLine{0}{Muss sich durchaus am steilsten Felsen ritzen.}
\VerseLine{1}{Das Köpfchen, kaum zum reinlich Denken gut,}
\Stanza
\VerseLine{0}{Will forschend in die tiefsten Gründe blitzen.}
\VerseLine{1}{Einmal im Leben fordert man ihr Blut:}
\VerseLine{0}{Sie wollen’s auf dem Schlachtfeld nur verspritzen!}
\end{BookVerse}
\PoemEnd
